% 初等整数估值与 LTE；独立知识节，不新增算法模板。
\section{素因子估值与升幂：先约去公共因子}\label{knowledge-lte}
\index{p-adic valuation}\index{lifting the exponent}\index{LTE}\index{升幂引理}
本节只处理整数的素因子重数。以下 $p$ 为素数，$a,b,n$ 为正整数；
对非零整数 $x$，定义 $v_p(x)=\max\{e\ge0:p^e\mid x\}$，负数按绝对值计算。
\textbf{不定义有限的 $v_p(0)$}：求差时先判断 $a=b$，此时 $a^n-b^n=0$，
被每个正整数整除，却没有最大的有限素因子重数，不能进入下面的估值公式或反复除 $p$ 的循环。

\subsection{基本运算与公共因子}
对非零整数 $x,y$，有
\[
v_p(xy)=v_p(x)+v_p(y),\qquad v_p(x^n)=n v_p(x).
\]
若 $y\mid x$，整数商满足 $v_p(x/y)=v_p(x)-v_p(y)$。
当 $x+y\ne0$ 时，$v_p(x+y)\ge\min(v_p(x),v_p(y))$；
两者估值不同才必定取等。证明只需提出较小的 $p$ 次幂：剩下两项中恰有一项不被 $p$ 整除；
估值相同时可能抵消，例如 $v_3(1)=v_3(2)=0$，但 $v_3(1+2)=1$。

计算 $a^n\pm b^n$ 时先记 $\alpha=v_p(a)$、$\beta=v_p(b)$。
若 $\alpha\ne\beta$，则直接得到
\[
v_p(a^n\pm b^n)=n\min(\alpha,\beta).
\]
若 $\alpha=\beta=t$，写成 $a=p^tA,b=p^tB$，其中 $p\nmid AB$，则
\[
v_p(a^n\pm b^n)=nt+v_p(A^n\pm B^n).
\]
差式仍要求 $a\ne b$。约去公共因子后，才能对 $A,B$ 检查下面的 LTE 条件；
不能把 $nt$ 误写为 $t$，也不能只凭 $p\mid a-b$ 就直接套公式。

\subsection{奇素数：差式与和式的条件不同}
以下两式均要求 $p$ 为\textbf{奇素数}且 $p\nmid ab$：
\[
\begin{aligned}
p\mid a-b,\ a\ne b
&\quad\Longrightarrow\quad v_p(a^n-b^n)=v_p(a-b)+v_p(n),\\
p\mid a+b,\ n\text{ 为奇数}
&\quad\Longrightarrow\quad v_p(a^n+b^n)=v_p(a+b)+v_p(n).
\end{aligned}
\]
差式不要求 $n$ 为奇数；和式不能漏掉这一条件。
这里 $p\nmid ab$ 只要求两数都不含 $p$，并不要求 $\gcd(a,b)=1$。

\textbf{为什么指数中每个 $p$ 恰好贡献一次。}
设 $n=mp^s$，$p\nmid m$。当 $p\mid a-b$ 时，
$(a^m-b^m)/(a-b)$ 模 $p$ 等于 $mb^{m-1}$，故指数的 $m$ 部分不增加估值。
再设 $x=y+h$，$p\nmid y$、$v_p(h)=r\ge1$。
二项式展开 $x^p-y^p$：一次项 $py^{p-1}h$ 的估值为 $r+1$，
其余项均至少为 $r+2$（末项用 $pr\ge r+2$）。
所以每次把指数乘 $p$ 都使估值增加 1，重复 $s$ 次即得差式。
和式可对 $a^2,b^2$ 使用差式，再用
$a^{2n}-b^{2n}=(a^n-b^n)(a^n+b^n)$：
当 $n$ 奇且 $a\equiv-b\pmod p$ 时，$p\nmid a-b$、$p\nmid a^n-b^n$，
两边相减只剩所求和式的估值。

\subsection{素数 2：先分指数奇偶}
若 $a,b$ 均为奇数且 $a\ne b$，则
\[
v_2(a^n-b^n)=
\begin{cases}
v_2(a-b),&n\text{ 为奇数},\\
v_2(a-b)+v_2(a+b)+v_2(n)-1,&n\text{ 为偶数}.
\end{cases}
\]
对奇数 $a,b$ 的和式（允许 $a=b$），有
\[
v_2(a^n+b^n)=
\begin{cases}
v_2(a+b),&n\text{ 为奇数},\\
1,&n\text{ 为偶数}.
\end{cases}
\]
奇指数时，差式与和式各自除以 $a-b$、$a+b$ 后的多项式均是奇数个奇数的和或交错和，
商为奇数。偶指数差式写 $n=2^sm$、$m$ 奇：先消去奇指数 $m$，再反复平方差分解；
首两因子是 $a-b,a+b$，以后每个 $a^{2^j}+b^{2^j}$（$j\ge1$）均模 4 余 2，恰好再贡献 $s-1$。
偶指数和式也因奇数的偶次幂模 4 余 1，估值恰为 1。
特别地，若 $4\mid a-b$，则 $v_2(a+b)=1$，差式可统一写成
$v_2(a^n-b^n)=v_2(a-b)+v_2(n)$；没有这个附加条件不能照搬奇素数公式。
若 $a,b$ 奇偶不同，和差均为奇数；若都偶，先按上一小节约去公共 2 因子。

\subsection{把大整数整除转成指数条件}
\textbf{例：求使 $2^6\cdot3^2\mid7^n-1$ 的最小正整数 $n$。}
奇数 $n$ 时 $v_2(7^n-1)=v_2(6)=1$，必不满足。
偶数 $n$ 时两种素数分别给出
\[
v_2(7^n-1)=1+3+v_2(n)-1=3+v_2(n),\qquad
v_3(7^n-1)=1+v_3(n).
\]
所以须且仅须 $v_2(n)\ge3$、$v_3(n)\ge1$，即 $24\mid n$，最小值为 24。
一般对非零整数 $X$ 和 $M=\prod p_i^{e_i}$，$M\mid X$ 等价于逐项检查
$v_{p_i}(X)\ge e_i$；无需真的构造可能远超机器字长的幂。
LTE 只能处理满足条件的局部因子，不能代替任意底数的乘法阶或离散对数问题。
实现时非零整数反复除 $p$ 即可求估值；$a+b$、差的绝对值及 $nt$ 的计算仍须防溢出。

\Needspace{120pt}
\textbf{条件漏一项，答案就可能错。}
\begin{itemize}
\item 不检查 $p\nmid ab$：$p=3,a=6,b=3,n=2$，实际 $v_3(27)=3$，
      硬套差式只得 $v_3(3)+v_3(2)=1$。
\item 不检查 $p\mid a-b$：$p=3,a=2,b=1,n=2$，实际 $v_3(3)=1$，
      右侧 $v_3(1)+v_3(2)=0$。
\item 将奇素数差式用于 2：$a=3,b=1,n=2$，实际 $v_2(8)=3$，
      错式只得 $v_2(2)+v_2(2)=2$。
\item 和式允许偶指数：$p=3,a=2,b=1,n=2$，实际 $v_3(5)=0$，
      错式却给 $v_3(3)+v_3(2)=1$。
\end{itemize}
定理条件参见
\href{https://oi-wiki.org/math/number-theory/lift-the-exponent/}{OI Wiki：升幂引理}。
本节限于整数估值、LTE 及整除建模，不代表完整的 $p$ 进数理论，也不新增实现或在线评测记录。

\textbf{与已有模板分工。}
扩展 Lucas 的阶乘赋值用 $v_p(N!)=\sum_{j\ge1}\lfloor N/p^j\rfloor$，
不是把整个阶乘硬套 LTE；见 \hyperref[compact-ExLucas]{ExLucas}
（第~\pageref{compact-ExLucas}~页，\ref{compact-ExLucas}~节）。
素数幂开根还要检查单位部分是否有根，单有估值条件并不充分；见
\hyperref[compact-PrimePowerRoots]{PrimePowerRoots}
（第~\pageref{compact-PrimePowerRoots}~页，\ref{compact-PrimePowerRoots}~节）。
